\subsection{可逆cross blockを持つtangent型の必要表示と全排除}
\label{subsec:h3-order2-22-tangent-normalform}
\begin{proposition}\label{prop:h3-order2-22-tangent-normalform}
半単純$(2,2)$・奇関係数2のAS4整域のtangent型で、偶cross block $K$が可逆とする。
この全領域は、$D(1+2D^2)\ne0$、$A=(1+D^2)/(2D)$の一変数必要表示
\begin{gather*}
 xu+uy-vx=0,\qquad xv+yu-ux=0,\\
 A(xy-yx)+y^2+u^2-v^2=0,\\
 Ax^2-yx+Dy^2-uv=0,\\
 -Ax^2-xy-Dy^2+vu=0,\\
 x^2+D(xy-yx)+u^2=0
\end{gather*}
で被覆される。ここで$x,y$は正、$u,v$は負生成元である。
この表示の全係数でAS4性やH3性を主張しない。
\end{proposition}
\begin{proof}
前節のtangent基底で$g=\left(\begin{smallmatrix}a&b\\c&d\end{smallmatrix}\right)$、
$ad-bc=1$とする。
混合topから偶cross blockは
\[
 K=\begin{pmatrix}-b&d&a&-c\\a&-c&0&0\\-d&0&c&0\\c&0&0&0\end{pmatrix},
 \qquad\det K=-c^4
\]
となる。従って$c\ne0$である。
$tangent$ pencilを保つ生成元変更は、その二基底に
\[
 E\mapsto E\begin{pmatrix}1&p\\0&r\end{pmatrix},\qquad
 F\mapsto F\begin{pmatrix}1&q\\0&r\end{pmatrix}
\]
という独立な$p,q$と共通の$r\ne0$の作用を誘導する。
実際、$x\mapsto\alpha x,y\mapsto\beta x+\delta y$、
$u\mapsto\gamma u,v\mapsto\eta u+\zeta v$で
$\delta/\alpha=\zeta/\gamma=r$とすれば、
$p=\beta/\alpha+\eta/\gamma,q=\beta/\alpha-\eta/\gamma$である。
そのgraph作用は$g\mapsto F_g gE_g^{-1}$であり、
$p=d/c,q=-a/c,r=1/c$により
$g=\left(\begin{smallmatrix}0&-1\\1&0\end{smallmatrix}\right)$へ正規化できる。

$K$可逆から純正quartic行列$P(A,B,C,D)$も可逆であり、$B\ne0$である。
相対homothetyとtop尺度で$B=1$とする。
$Q=K^tP^{-1}K$の四cyclic条件を分母を払って書くと、
$X=-Cc+D(a-d)$に対して、そのうち二つは
\[
 2BDcX=0,\qquad -2B(Bc+D(a+d))X=0.
\]
$Bc\ne0$なので$X=0$が従う。
残る条件は
\[
 -4ADc^2+2B^2c^2+C^2c^2+2D^2(ad-bc)=0.
\]
正規化した$g$では$C=0$、$4AD=2+2D^2$となり、$D\ne0$である。
さらに$\det P=-(1+2D^2)$なので表示した開条件が必要となる。
奇graph二式と偶blockの先頭四行$[P\ K]$を並べれば、主張した六関係を得る。
\end{proof}

\texttt{verify\_cubic\_order2\_22\_tangent\_normalform.py}は
全係数のcross行列式、四Schur条件、pencil安定化作用と一変数の行列式を厳密検算する。
$K$が退化する$c=0$境界は別に残り、この命題の正規化で除外しない。
指定正常三次は次定理で排除する。任意係数でのAS4存在は主張しない。

\begin{theorem}[tangent可逆部分の全値排除]
\label{thm:h3-order2-22-tangent-fullrank-exclusion}
半単純$(2,2)$・奇関係数2のtangent型で、偶cross blockが可逆なH3対象は存在しない。
より強く、前命題の表示が指定中山作用を持つAS4となる場合、
正符号三次成分で$x,y$の双方と可換する元は零である。
これは$D(1+2D^2)\ne0$の全値を含む。
\end{theorem}
\begin{proof}
正符号の三次自由語空間を$V_3^+$、四次自由語空間を$V_4^+$とする。
その次元は32と128である。
明示的な33個の多項式線型汎関数
$\lambda_i:V_4^+\to\mathbb Q[D]$を入力証明書に取る。
各汎関数が六つの二次関係から生じる全偶四次関係を消すことは、
4,752本の多項式恒等式で直接検証できる。
従って任意の$D\ne0$で$B_4^+$上の線型汎関数を与える。
三つの関係を$2D$倍して検証しているが、これはこの開集合上で可逆な操作である。

これらを二つの交換子へ適用して
\[
 N:V_3^+\longrightarrow k^{66},\qquad
 F\longmapsto\bigl(\lambda_i(Fx-xF),\lambda_i(Fy-yF)\bigr)_i
\]
を作る。係数は全て$\mathbb Q[D]$内にある。
この$66\times32$行列の、入力証明書に指定した九つの$10\times10$ minorの
最大公約因子は非零有理数倍を除き$D^{12}$である。
従って任意の$D\ne0$で少なくとも一つのminorが非零となり、
$\operatorname{rank}N\ge10$である。
有限標本からの推論ではなく、多項式minorの正確な最大公約因子による全値の被覆である。

指定AS4を仮定すると、命題\ref{prop:h3-order2-mixed-rank}の等変Hilbert級数から
$\operatorname{Tr}(\mu_B|_{B_3})=0$である。
$\dim B_3=20$かつ中山が半単純なので$\dim B_3^+=10$、
従って三次関係空間$\mathcal I_3^+\subset V_3^+$の次元は22である。
交換子は四次関係へ入るため$\mathcal I_3^+\subset\ker N$である。
上の階数から$\dim\ker N\le22$なので、両者は一致する。
従って$F\in B_3^+$が$x,y$と可換すれば、もともと関係であり$B$内では零である。
H3に必要な$G$は$\mu_B(G)=G$で、$x,y$と可換するので、この矛盾により排除される。
\end{proof}

汎関数は、一般点の19個のannihilatorへ
$h(D)=10D^6-D^4-4D^2-1$の零点で不足する14個を補って構成した。
ただし証明はこの構成法を仮定せず、保存した33個そのものの全恒等式を検査する。
\texttt{cubic\_order2\_22\_tangent\_annihilator\_input.json}が全係数と九minorの位置を持ち、
\texttt{verify\_cubic\_order2\_22\_tangent\_exclusion.py}が恒等式・行列式・最大公約因子を再生する。
有理数の定数分母も入力に保持し、多項式化による係数の脱落を防いでいる。
この結果の後もtangentの退化cross境界$c=0$、secant型、非半単純型は残る。

\begin{theorem}[退化境界と全tangent型の排除]
\label{thm:h3-order2-22-tangent-total-exclusion}
半単純$(2,2)$・奇関係数2のtangent型は、退化cross境界も含めて全てH3から排除される。
従ってこの領域の残る候補はsecant型だけである。
\end{theorem}
\begin{proof}
可逆部分は前定理で済んだので、$c=0$、$ad=1$とする。
純正topを$P(A,B,C,D)$と書く。
$y^2$行と、$d$倍の$xy$行に$a$倍の$yx$行を加えた行は、純正二次関係であり、
その係数は
\[
 q=(B,D,-D,0),\quad
 r=(A(d-a),dC-aB,-dB-aC,D(d-a))
\]
である。整域性からそのspanは高々一次元で、非零rank-one tensorを含めない。
$q\ne0$なら$D\ne0$である。二つのminor
$D^2(d-a)$と、$d=a$での$-2aBD$から$d=a=\epsilon=\pm1,B=0$となる。
$v$を$u$で移動してgraphの上右係数$b$を消せる。

まず$q=0$、すなわち$B=D=0$の場合を排除する。
同じ移動で$b=0$とする。
$C=0$なら$r=A(d-a)x^2$なので$A(d-a)=0$が必要である。
$A=0$でも$d=a$でも、偶関係空間の$U^2$射影は高々二次元となる。
従って四次元の偶関係の核は$V^2$内に二次元以上となり、整域性に反する。
よって$C\ne0$である。
$xy$行から
\[
 u^2=-a^{-1}x(Ax+Cy),\qquad ux=-a^{-1}xu
\]
を得る。$uu^2=u^2u$へ代入し、整域で$x$を消去すると
\[
 uy+a yu+(A/C)(a-a^{-1})xu=0
\]
となる。これは$v$を含まない新しい奇二次関係である。
元の奇関係は$xu+a ux$と、$xv$係数1を持つtangent関係の二本だけなので、
この新関係はそのspanに入らない。これは$B_2$の関係数に反する。
従って$q=0$の場合はAS4整域自体が存在しない。

残る$D\ne0$、$g=\epsilon I$を扱う。
$\epsilon=1$で偶middle blockの添字$(yy,xy,xx,uv)$の主minorは$D^2$である。
この四行に対するSchur条件は、純負cyclic block
\[
 Q=\begin{pmatrix}
 a_0&b_0&b_0&c_0\\b_0&d_0&c_0&e_0\\b_0&c_0&d_0&e_0\\c_0&e_0&e_0&f_0
 \end{pmatrix}
\]
について$a_0=c_0=d_0=f_0=0$と同値である。
$\epsilon=-1$ではcross blockの符号を変える合同変換となるので同じ条件である。
$e_0=0$なら$V^2$射影が高々二次元となり、再び純正関係が二次元以上となる。
従って$e_0\ne0$であり、$v^2$行から$uv+vu=0$を得る。
既知の純正関係から$xy=yx$である。
残る偶関係は
\[
 u^2=-\epsilon(Ax^2+Cxy+Dy^2),\qquad
 v^2=-\epsilon e_0^{-1}x^2-(b_0/e_0)u^2.
\]
$y$の$x$移動で$C=0$にできる。この移動は$x u+\epsilon ux=0$により奇関係を変えない。
従って、ある係数$\alpha,\delta,\beta,\gamma$により
\begin{gather*}
 xy=yx,\quad xu=-\epsilon ux,\quad uv=-vu,\\
 xv-\epsilon vx=-(yu+\epsilon uy),\\
 u^2=\alpha x^2+\delta y^2,\quad
 v^2=\beta x^2+\gamma y^2,\qquad
 \delta(\alpha\gamma-\beta\delta)\ne0
\end{gather*}
という必要表示になる。

この表示では四生成元の平方が全て中心である。
実際、$T=yu+\epsilon uy$について$xT=-\epsilon Tx$なので、奇関係から$x^2v=vx^2$となる。
これで$x^2$は中心である。
$u^2$は正可換平面の二次式であり、$uv=-vu$も使えば全生成元と可換する。
$\delta\ne0$から$y^2$も中心となり、最後に$v^2$も中心となる。

二つの非零二次式
\[
 W=yu+\epsilon uy=-(xv-\epsilon vx),\qquad N=yv-\epsilon vy
\]
を取る。$N$の二語は元の奇関係に現れず、$W$は$xv$語を消すことなく奇関係に入れないので、
どちらも$B_2$で非零である。
中心平方と上の関係から、その正規化作用は
\[
 \sigma_W=(-\epsilon,\epsilon,\epsilon,-\epsilon),\qquad
 \sigma_N=(\epsilon,-\epsilon,\epsilon,-\epsilon)
\]
となる。例えば$W$の$y,u$に対する式は中心平方から、$x,v$に対する式はもう一つの表示から従い、
$N$の$x,u$に対する式はこれらを代入して得る。
従って$M=WN\ne0$の正規化作用は$-\mu_B$である。
一方$\mu_B(W)=-W,\mu_B(N)=-N$なので$\mu_B(M)=M$である。
H3の$G$は$GM=MG$を満たすが、$M$の正規化式と$\mu_B(G)=G$、$\deg G=3$から
$MG=-GM$でもある。これは整域性と標数0に反する。
\end{proof}

\texttt{verify\_cubic\_order2\_22\_tangent\_boundary.py}は
境界の二minor、全係数のSchur条件、および両符号の二つの正規二次式の全共変式を厳密検証する。
最後の共変式は、中心平方と必要な三二次関係による三次自由語空間の行恒等式として再生する。
今回の境界排除は独立な証明であり、as4から新しい被覆定理を輸入していない。
